%% ================================================================================================
%%  6. The procedure: Condorcet with a blocking quota
%%  Notes: research-agenda.md §6; proofs in A_proofs.tex; algorithms in B_algorithms.tex
%%  Code: simulations/python/voting.py (qvc, --check), examples.py (table examples, --check)
%% ================================================================================================
\section{The procedure: Condorcet with a blocking quota}\label{sec:procedure}

\subsection{Design principles}

Sections~\ref{sec:limits}--\ref{sec:aggregation} translate into six principles.
\begin{enumerate}[label=P\arabic*]
  \item \textbf{One simultaneous ballot on the whole agenda.} Voters judge all options at once (H2): no
  sequential or amendment agenda, hence no power for whoever orders the votes (Section~\ref{sec:veto}).
  \item \textbf{Ranking and blocking are separate acts.} The ballot has a ranking and an acceptability
  line; only the line is a veto. This is the separation the literature favours
  \citep{BoutonLlorenteSaguerMalherbe2018}.
  \item \textbf{Vetoes are held by coalitions, not individuals.} An option is excluded only if a blocking
  minority, a share $\theta$ of the voters, finds it unacceptable: minority protection without holdouts
  (Propositions~\ref{prop:majority-veto} and~\ref{prop:gridlock}).
  \item \textbf{The status quo is an option like the others.} It can be ranked, vetoed and elected, but
  it is never the default: vetoes cannot become a vote for inertia.
  \item \textbf{The aggregation is clone-independent.} The open agenda (H1) must not reward flooding
  (Section~\ref{sec:clones}).
  \item \textbf{A decision is guaranteed.} If every option is vetoed by a blocking minority, the
  least-vetoed options are kept rather than none (R2).
\end{enumerate}

\subsection{The procedure}

\begin{definition}[Admissible options; $\theta$-compromise]\label{def:compromise}
Let $0 < \theta \le 1/2$. The \emph{admissible} options are
$\Adm = \{x \in \Opts : \veto(x) < \theta n\}$ if this set is not empty, and the least-vetoed options
$\arg\min_x \veto(x)$ otherwise. The \emph{$\theta$-compromise} is the Condorcet winner of $\Adm$, when it
exists.
\end{definition}

The definition is lexicographic: first, the outcome must be acceptable to all but a minority smaller than
$\theta$, if any option is; second, among such options, the majority decides. It is our reading of R1.
When the Condorcet winner of the whole agenda is acceptable to all but such a minority, it is the
$\theta$-compromise; when a blocking minority finds it unacceptable, as in Example~E1, the compromise is
the best option, in the majority sense, that the minority can live with.

\begin{algorithm}[htbp]
\caption{\QVC($\theta$): Condorcet with a blocking quota}\label{alg:qvc}
\begin{algorithmic}[1]
\Require options $\Opts$ (all proposals and the status quo), ballots $(\wpref_i, V_i)_{i \in \Voters}$,
  quota $\theta \in (0, 1/2]$, an order of the options drawn at random before the vote (for exact ties)
\State $\veto(x) \gets |\{i : x \in V_i\}|$ for every $x \in \Opts$ \Comment{veto counts}
\State $\Adm \gets \{x : \veto(x) < \theta n\}$
\If{$\Adm = \emptyset$} $\Adm \gets \arg\min_x \veto(x)$ \Comment{fallback: never empty}
\EndIf
\State $\Maj(x,y) \gets |\{i : x \pref_i y\}|$ for $x, y \in \Adm$ \Comment{pairwise counts}
\State $W \gets$ Schulze winners of $\Adm$ \Comment{Appendix~\ref{app:algorithms}}
\If{$|W| > 1$} $W \gets \arg\min_{x \in W} \veto(x)$, then the Schulze winners of $W$, then the first
  in the random order
\EndIf
\State \Return the option in $W$; publish $\veto$, $\Adm$ and the matrix $\Maj$
\end{algorithmic}
\end{algorithm}

Algorithm~\ref{alg:qvc} computes the $\theta$-compromise when it exists and completes it with the
Schulze method otherwise. Ranked Pairs could replace Schulze without changing any of the properties below.
The computation takes $O(nm^2 + m^3)$ operations, a fraction of a second for hundreds of options.

\subsection{Properties}

\begin{proposition}[Decisive, with no default advantage]\label{prop:decisive}
\QVC($\theta$) selects exactly one option, in one ballot round, for every profile. It treats every option,
the status quo included, in the same way, and its outcome does not depend on the order in which the
options were proposed (except through the random order that breaks exact ties).
\end{proposition}

\begin{proposition}[Acceptability guarantee]\label{prop:acceptability}
If some option is vetoed by fewer than $\theta n$ voters, so is the winner of \QVC($\theta$): more than a
share $1-\theta$ of the voters find it acceptable. Otherwise the winner is among the least-vetoed
options. Hence the winner is vetoed by at least $\theta n$ voters only if every option is, and no rule
has a lower probability of electing an option vetoed by a blocking minority.
\end{proposition}

\begin{proposition}[Majority within the admissible options]\label{prop:condorcet-admissible}
The winner of \QVC($\theta$) is in the Smith set of $\Adm$ and is not Pareto-dominated. If $\Adm$ has a
Condorcet winner, it wins; in particular the Condorcet winner of the whole agenda wins whenever fewer
than $\theta n$ voters veto it.
\end{proposition}

\begin{proposition}[Independence of clones]\label{prop:qvc-clones}
\QVC($\theta$) is independent of clones (Definition~\ref{def:clones}), up to the random order that breaks
exact ties.
\end{proposition}

The proofs are in Appendix~\ref{app:proofs}; the code checks Propositions~\ref{prop:decisive}--\ref{prop:qvc-clones} on random profiles
(Appendix~\ref{app:simulation}). The guarantees have a price, which is the subject of the
last result.

\begin{proposition}[Strategic blocking]\label{prop:blocking}
Under \QVC($\theta$), any coalition of at least $\theta n$ voters can make any option inadmissible by
vetoing it, which defeats that option whenever another option stays admissible. If $\theta \le (m-1)/m$,
and in particular for every $\theta \le 1/2$, there are profiles in which, whatever the winner $w$, at
least $\theta n$ voters prefer some other option to $w$, and can thus block $w$ by vetoing it although
they find it acceptable.
\end{proposition}

Proposition~\ref{prop:blocking} is the strategic face of Proposition~\ref{prop:majority-veto}: the power to
block an unacceptable option is also the power to block an acceptable one. \citet{Greenberg1979} gives
the general threshold: an option that no coalition of size $\theta n$ wants to replace exists for every
profile only if $\theta > (m-1)/m$, or $\theta > d/(d+1)$ in $d$ dimensions, that is, only if the veto is
weaker than majority rule. A sincere veto is cheap to cast and a strategic one is no dearer, so
safeguards must come from the procedure around the vote (Section~\ref{sec:safeguards}). The proportional
veto core avoids the problem, since it is defined as the set of options no coalition can block given its
proportional veto power \citep{Moulin1981,Moulin1982}, but it fails independence of clones
(Proposition~\ref{prop:pvc-clones}), which H1 makes essential.

\subsection{Worked examples}\label{sec:examples}

Table~\ref{tab:examples} applies nine rules to four small profiles.

% Generated by simulations/python/examples.py -- do not edit by hand
\begin{table}[htbp]\centering\footnotesize
\setlength{\tabcolsep}{4pt}
\caption{Winners in the worked examples}\label{tab:examples}
\sbox{\tabbox}{\begin{tabular}{lrrrrrrrrr}
\toprule
Example & Plur. & Borda & Schul. & Fallb. & PVC & Unan. & Q 25\% & Q 33\% & Q 50\% \\
\midrule
E1 (compromise) & A & B & A & B & B & B & B & B & A \\
E2 (flooding) & A & A & A & A & A & B & B & B & A \\
E3 (cycle) & A & A & A & B & A & C (stuck) & B (fb) & A & A \\
E4 (status quo) & X & X & X & Y & X & S (stuck) & X & X & X \\
\bottomrule
\end{tabular}}
\usebox{\tabbox}
\tabnotes{E1: 51 voters rank A, B, C and veto C; 49 rank B, C, A and veto A. E2: E1 after the majority proposes two variants of A, ranked just below A by the majority and last by the minority, which vetoes them. E3: 35 voters rank A, B, C; 33 B, C, A; 32 C, A, B; each bloc vetoes its last option. E4: 99 voters rank X, Y, S and veto S; one voter ranks S, Y, X and vetoes X and Y. The status quo is C in E1--E3 and S in E4. Stuck: no proposal passes and the status quo stays; fb: every option is vetoed by the quota and QVC keeps the least-vetoed ones. Fallb.: fallback bargaining; PVC: Schulze within the proportional veto core; Unan.: unanimity with the status quo as default; Q $x$: QVC with a blocking quota $\theta = x$.}
{Profiles of Section~\ref{sec:examples}; \texttt{examples.py}.}
\end{table}


\begin{itemize}
  \item \textbf{E1, majority or compromise.} $A$ is the Condorcet winner (51 to 49 against both others)
  and plurality and Schulze elect it. The 49\% minority, however, ranks $A$ last and vetoes it, while $B$
  is everyone's first or second choice. Borda's rule, fallback bargaining, the proportional veto core and
  \QVC\ with $\theta \le 49\%$ elect $B$; \QVC(1/2) treats a 49\% veto as insufficient and elects $A$.
  Which answer is right is the choice of $\theta$.
  \item \textbf{E2, flooding.} The majority adds two variants of $A$. Borda's rule, fallback bargaining
  and the proportional veto core now elect $A$ (Proposition~\ref{prop:pvc-clones}); \QVC\ is unchanged,
  because each variant inherits the minority's vetoes.
  \item \textbf{E3, a cycle.} $A$ beats $B$ (67--33), $B$ beats $C$ (68--32), $C$ beats $A$ (65--35);
  Schulze breaks the weakest link and elects $A$. With $\theta = 1/4$ every option is vetoed by about a
  third of the voters, so \QVC\ falls back on the least-vetoed option, $B$; with $\theta = 1/3$, $C$ (35
  vetoes) is excluded and $A$ wins. Unanimity keeps the status quo $C$, since neither proposal is
  preferred to it by everyone.
  \item \textbf{E4, the status-quo trap.} One voter prefers the status quo to everything and vetoes both
  proposals. Unanimity keeps the status quo against 99 voters; \QVC\ elects $X$, since one veto is far
  below any blocking quota.
\end{itemize}

\subsection{Choosing the quota and the details}\label{sec:theta}

\paragraph{The quota.} $\theta$ sets the size of the smallest minority whose objection counts. The
simulations of Section~\ref{sec:simulation} (Figure~\ref{fig:theta}) show three regimes.
\begin{itemize}
  \item \emph{At 15--20\% and below}, in an assembly of a hundred voters where each voter vetoes 35--40\% of
  the options, every option is vetoed by the quota and the fallback decides in most elections: the
  procedure elects the least-vetoed option, which is approval voting on acceptability. Each individual
  veto is then only a negative vote, and the result is easily manipulated.
  \item \emph{Between 25\% and 40\%}, the veto excludes the options that a substantial minority rejects
  and the majority decides among the others; the fallback is rare.
  \item \emph{From 40\% upwards}, the veto almost never binds before majority rule would and \QVC\ is the
  Schulze method.
\end{itemize}
We recommend $\theta = 1/3$ as the default: it is the blocking third that corresponds to the two-thirds
supermajorities common in statutes, close to the 35\% blocking minority of the Council of the European
Union and to the 36\% implied by the 64\% rule of \citet{CaplinNalebuff1988}. A group that values
consensus more than majority can go down to $\theta = 1/4$; going lower brings back the problems of the
individual veto. In a small group, the quota is a number of voters, $\lceil \theta n \rceil$: with ten
voters and $\theta = 1/3$, four vetoes exclude an option. Below five voters, a quota of two vetoes keeps a
single member from blocking alone.

\paragraph{The status quo.} By default the status quo is an ordinary option (P4). Where a change is
costly or irreversible, the group may wish to protect it, as Debian does: the winner must then beat the
status quo by a stated ratio (2:1, say), failing which the status quo stays. This reintroduces inertia
on purpose and should be reserved for such decisions.

\paragraph{The ballot.} A ranking with ties allowed, plus the acceptability line. With many options,
a grade-style form is easier: each option receives one of ``excellent'', ``good'', ``fair'', ``poor'' or
``unacceptable'', voters may order options within a grade, and the ranking and the vetoes are read from
the form. Options left blank are ranked last and not vetoed. Results should publish the veto counts, the
admissible set and the pairwise matrix, so that any member can check the outcome.

\subsection{Safeguards against strategic vetoing}\label{sec:safeguards}

Proposition~\ref{prop:blocking} cannot be removed by the counting rule; it can be contained by the
procedure around it.
\begin{itemize}
  \item \textbf{Named and reasoned vetoes.} A veto is public and states a reason, as an objection must in
  the IETF's rough consensus \citep{Resnick2014rfc7282}. Responsibility for one's vote changes behaviour in
  committees \citep{BreitmoserValasek2024}, and social constraints are what keeps formal vetoes in check
  in consensus bodies \citep{Krick2017}. A chair or a panel may rule a veto without a reason void.
  \item \textbf{One revision round, only when needed.} If the fallback was used (no option acceptable
  to more than $1-\theta$ of the voters), the results are published, the nomination opens again for a
  fixed time for compromise proposals, and a second ballot is held; its result is final. Two rounds at
  most keep R2 while giving a divided group a chance to find a compromise, in line with the evidence that
  one or two rounds suffice and more rounds strengthen vetoers
  \citep{GersbachValkanova2024,Nunnari2021}.
  \item \textbf{Veto budgets across decisions, not within one.} In a body that decides repeatedly, a
  limited stock of vetoes per member and per period, which can be saved for decisions that matter more,
  discourages routine blocking \citep[on storable votes]{Casella2005}. A budget within one decision would
  break independence of clones: each clone added would cost a veto.
  \item \textbf{No lower quota than needed.} The manipulability of the procedure falls quickly as
  $\theta$ rises (Figure~\ref{fig:theta}).
\end{itemize}
